\section{Methods}

\subsection{Frozen representation pipeline}

We used the \texttt{brain-bzh/reve-base} checkpoint with encoder parameters
frozen, evaluation mode enabled, and inference-mode extraction. EEG was mapped
to the ordered EGI HydroCel E1--E128 layout without spatial interpolation,
resampled to 200 Hz, band-pass filtered from 0.5 to 99.5 Hz, standardized and
clamped, and divided into non-overlapping two-second windows. At most 120
seconds were retained per subject. Windows were never allowed to cross MIPDB
condition boundaries. Representations from the last two transformer layers were
materialized once and content-addressed; all heads consumed these immutable
caches.

\subsection{Probe heads}

The baseline averages final-layer tokens and applies one affine regressor. The
penultimate-layer linear probe makes the same readout from the preceding layer,
isolating layer choice without adding capacity. The rich-statistics residual
head adds a zero-initialized correction over per-feature standard deviation,
range, mean absolute deviation, and mean absolute magnitude. The multi-query
head uses two deterministic signed-basis attention queries and a
zero-initialized residual over weighted location, spread, range, deviation, and
magnitude statistics. Zero initialization makes both residual heads exactly
match the baseline prediction at initialization, so subsequent differences
arise through training rather than an unmatched initial function.

\subsection{Training and selection}

Each head was trained for seeds 33--42 using global seeded window shuffling,
AdamW with learning rate 10 to the minus 4 and weight decay 0.05, batch size 64,
OneCycleLR with cosine annealing, gradient clipping at norm 1, and mean-squared
error. Training had a maximum of 40 epochs and patience 7. For each seed and
head, the earliest checkpoint attaining the best HBN validation Pearson
correlation was selected. Neither MIPDB targets nor predictions influenced
training, checkpoint selection, candidate membership, or thresholds.

\subsection{External estimand and descriptive metrics}

For each candidate head and optimization seed, we computed subject-level Pearson
correlation on the primary external cohort. The primary seed-level contrast is
the candidate Pearson correlation minus the matched baseline correlation, and
the estimand is its arithmetic mean over the ten matched seeds. We also report
MAE, RMSE, R-squared, calibration
slope and intercept, the sample standard deviation across seed correlations,
the number of positive, zero, and negative seed contrasts, and the worst seed
contrast. These additional metrics are descriptive and do not replace the
predeclared decision rule.

\subsection{Uncertainty and decision rule}

The interval estimate uses 10,000 deterministic hierarchical paired bootstrap
iterations. Each iteration resamples seeds with replacement and subjects with
replacement, preserving candidate--baseline pairing. An exact one-sided
seed-randomization test enumerates all 1,024 sign assignments for the ten
paired seed deltas. Raw p-values for the three candidates are adjusted in a
fixed order with Holm's step-down method.

A candidate establishes stable external improvement only if all conditions are
true: the primary cohort meets the 50-subject minimum; adjusted p is below
0.05; the
95\% bootstrap lower bound is above zero; at least eight of ten seeds improve;
and the worst seed delta is at least minus 0.01. This conjunction was fixed before
external analysis. Failure of any condition blocks a superiority claim.
